\section{Discussion and limitations}\label{sec:discussion}

\subsection{Interpretation of the results}\label{sec:interpretation}

The two kinds of round differ in the supervision they give. A FedAvg SFT
round raises the likelihood of one gold SQL query per question, the same
query in every round, whereas a preference round compares queries of the
current global model that the client has executed, so its training
signal changes with the model. FedAvg SFT does not improve after its
third round on either benchmark, while from its model after
\setting{rounds.warmup} rounds a preference round adds more than a
further FedAvg SFT round, significantly so on Spider test
(Section~\ref{sec:main}). We read this as the
gold SQL queries being exhausted as a training signal sooner than the
executed queries of the model.

The preference rounds reduce execution errors, not incorrect results
(Section~\ref{sec:analysis}), and the composition of the pairs is
consistent with this. The incorrect sample is an execution error in
\fact{pairs.bird.pref1.neg_error_pct}\% of the pairs of the first
preference round on BIRD and \fact{pairs.spider.pref1.neg_error_pct}\%
on Spider, and in \fact{pairs.bird.pref2.neg_error_pct}\% and
\fact{pairs.spider.pref2.neg_error_pct}\% in the second. We attribute
the difference between the two kinds of error to what a pair exposes. A
query that fails to execute differs from the correct sample in a way
that the pair shows directly, such as a name that does not exist in the
schema or a malformed clause, whereas a query that returns another
result can differ from it in its reading of the question.

Most of the reduction comes from the comparison of the two samples and
not from training on the correct samples of the model. Without the
preference term, execution errors are more numerous after one preference
round on both benchmarks and after two on BIRD dev, where the variant
ends with \nerr{bird}{rft}{4} against \nerr{bird}{fedavg}{4} for FedAvg
SFT and \nerr{bird}{k2}{4} for \method{}
(Section~\ref{sec:ablation}). The same variant separates two sources of
the EX gain on BIRD dev, both of which use the execution labels: training
on the samples labeled correct puts it \gap{bird}{fedavg}{4}{rft}{4}
points above FedAvg SFT, and the preference term adds
\gap{bird}{k2}{4}{rft}{4}. The outcome of each question bounds what the reduction is worth. Of the
execution errors of FedAvg SFT, more return an incorrect result with
\method{} than become correct on BIRD dev (\fact{moved.bird.error.wrong}
against \fact{moved.bird.error.correct}), and fewer on Spider test
(\fact{moved.spider.error.wrong} against
\fact{moved.spider.error.correct}). On both benchmarks, correct
predictions rise by less than execution errors fall. Removing an
execution error is not the same as answering the question correctly.

On Spider test, EX does not differ significantly across the numbers of
warm-up rounds tested, but without warm-up the EX gain over FedAvg SFT does not carry to TS on
the three sets with a suite, and with \setting{rounds.warmup} warm-up
rounds it does (Section~\ref{sec:perturbed}). An execution label, like
EX, is a match of results on one database. The suites show that
preference rounds started from the pretrained model yield more
predictions that match the gold result on the evaluation database and
return another result on its content variants. The experiments do not
show why.

Renewing the pairs and the reference model, like the second preference
round itself, adds accuracy on BIRD dev and not on Spider test
(Section~\ref{sec:ablation}). This is consistent with the smaller share
of training questions that yield pairs on Spider, where fewer questions
lack a correct sample after warm-up (Section~\ref{sec:warmup}). From the
first to the second preference round, the share falls from
\fact{cost.bird.pref1.paired_pct}\% to \fact{cost.bird.pref2.paired_pct}\%
on BIRD and from \fact{cost.spider.pref1.paired_pct}\% to
\fact{cost.spider.pref2.paired_pct}\% on Spider.

\subsection{Scope and limitations}\label{sec:limitations}

A client needs gold SQL queries for its training questions and databases
on which it can execute queries, and questions without a gold SQL query
are not used. A question on a database without rows, or one whose gold
SQL query fails to execute or exceeds the time limit, yields no pair,
and the number of questions of the last kind differs between runs
(Section~\ref{sec:datasets}).

An execution label is a match of results on one database, so an
incorrect query can be labeled correct when its result matches by
chance. TS is the check that the study has: the official suite of
the Spider dev databases for Spider dev, Spider-Syn and Spider-Realistic,
and a generated suite on BIRD dev, whose agreement with an official suite
could be measured on Spider dev only. On BIRD dev, this suite rejects
about a third of the EX-correct predictions that read other tables or
columns than the gold SQL query, for FedAvg SFT and for \method{} alike
(Section~\ref{sec:analysis}), which gives the size of the effect at
evaluation. We have no such measure for the labels of the training
samples. Spider test and Spider-DK are judged by EX alone.

Each model comes from one training run, with one seed and one draw of
the partition, so the tests cover the variation over evaluation
questions and not the variation between runs.

\method{} needs more client computation over the rounds: in
\setting{rounds.total} rounds, \fact{cost.bird.k2.gpu_total} against
\fact{cost.bird.fedavg4.gpu_total} GPU minutes for FedAvg SFT on BIRD
and \fact{cost.spider.k2.gpu_total} against
\fact{cost.spider.fedavg4.gpu_total} on Spider, plus the labeling time
(Section~\ref{sec:cost}). The models are compared at equal numbers of
rounds, that is, at equal communication, and only on Spider does a
comparison give the baseline more GPU time, through its fifth round. The
preference rounds also use a lower learning rate and weight decay than
the FedAvg SFT rounds (Section~\ref{sec:implementation}). No FedAvg SFT
run uses these settings, and the variant without the preference term,
which shares them, is the only control for them.

The setups have \fact{split.bird.a0.1.k5.clients} clients, or
\fact{split.spider.a0.5.k10.clients} in one setup, and all of them take
part in every round. Partial participation and federations of many
clients are not studied. With \setting{rounds.pref} preference rounds,
the pairs and the reference model are renewed once, and the ablation
study does not separate the new pairs from the new reference model. The
number of warm-up rounds is a setting that has to be chosen.

Privacy in this study is data isolation: no differential privacy is
applied, and what an adapter reveals about the data of a client is not
analyzed. The pair count that a client sends depends on how the model
performs on its data.

The models have \setting{models.size.min} to \setting{model.size}
parameters and are trained with LoRA, and the benchmarks consist of
SQLite databases with questions in English. With the smaller model, the
gain over FedAvg SFT is smaller (Section~\ref{sec:breadth}).
